
\section{Appendix A. Mathematical Derivations}

\subsection{A.1 SVD Truncation Error}

The rank-$r$ SVD approximation satisfies:

\[
\|W - W_r\|_F = \sqrt{\sum_{i > r} \sigma_i^2}
\tag{A.1}
\]

This is optimal (Eckart-Young).

\subsection{A.2 SALT Parameterization}

We define:

\[
W_\text{SALT}(\gamma, A, B) = U_r \text{diag}(\Sigma_r \odot \gamma) V_r + AB
\tag{A.2}
\]

\subsection{A.3 SALT at Initialization}

At init:

\begin{itemize}
\item $\gamma = \mathbf{1}$
\item $B = 0$

\end{itemize}
So:

\[
W_\text{SALT}(\mathbf{1}, A, 0) = U_r \Sigma_r V_r + A \cdot 0 = W_r
\tag{A.3}
\]

The model starts as the truncated SVD.

\subsection{A.4 Expressiveness}

\textbf{Claim.} The family of matrices representable by $W_\text{SALT}$ includes 
all rank-$r$ matrices, and additionally any matrix of the form 
$W_r + AB$ where $\text{rank}(AB) \leq r$.

\textbf{Proof.} Setting $\gamma$ and $A, B$ appropriately, we can reach any 
matrix in $\{W_r + AB : \text{rank}(AB) \leq r\}$. $\square$

\textbf{Effective rank:} up to $2r$.

\subsection{A.5 Parameter Count}

Let $W \in \mathbb{R}^{m \times n}$.

\begin{table}[H]
\centering
\small
\begin{tabular}{ll}
\toprule
\textbf{Component} & \textbf{Params} \\
\midrule
$U_r$ & $m \cdot r$ \\
$\Sigma_r$ & $r$ \\
$V_r$ & $r \cdot n$ \\
$\gamma$ & $r$ \\
$A$ & $m \cdot r$ \\
$B$ & $r \cdot n$ \\
bias & $m$ \\
\textbf{Total} & $2r(m + n) + 2r + m$ \\
\bottomrule
\end{tabular}
\end{table}

For $m = n = 768$, $r = 16$:

\[
\text{Total} = 2 \cdot 16 \cdot 1536 + 32 + 768 = 49,184
\tag{A.4}
\]

Full W = 589,824. Ratio = 11.99$$\times$.

\subsection{A.6 Inference Complexity}

\subsubsection{Naive Forward}

\[
T_\text{naive} = T_\text{reconstruct} + T_\text{matmul} = \mathcal{O}(d_\text{out} r d_\text{in}) + \mathcal{O}(BL d_\text{out} d_\text{in})
\tag{A.5}
\]

\subsubsection{Fused Forward}

\[
T_\text{fused} = \mathcal{O}(BL r (d_\text{out} + d_\text{in}))
\tag{A.6}
\]

\textbf{Ratio:} $T_\text{naive} / T_\text{fused} \approx \frac{d_\text{in} + d_\text{out}}{r (1 + 1/BL)}$. 
For $B = 8, L = 128, r = 16$: ~100$$\times$ more FLOPs in the naive version (mostly 
from the matmul, not reconstruction).

\subsection{A.7 GQA Pooling}

For $H_{kv} < H$, the source has more KV heads. We pool:

\[
W_{k,\text{pooled}}[i, :] = \frac{1}{\text{rep}} \sum_{j=1}^{\text{rep}} W_{k}[i \cdot \text{rep} + j, :]
\tag{A.7}
\]

This is a mean-pooling over groups of size $\text{rep} = H / H_{kv}$.
\begin{lstlisting}

---
\end{lstlisting}